\chapter{Conclusion}
\label{ch:conclusion}

\openingpara
{This dissertation advances affective Human--Building Interaction} {as an orientation for specifying and investigating the situated complexity of indoor experience and bringing this understanding into the design of future interactive indoor environments. It maps the experiential agenda of research on interactive indoor environments, examines how environmental conditions acquire significance through affect, develops multisensory biophilic design knowledge, and investigates one sound-based instantiation through an AI-transformed \sidenote{Throughout this dissertation, I use the term \emph{transformed} to refer to AI-transformed audio. The RAVE model receives an incoming audio signal and generates a corresponding output through the learned nature-sound representation.} office soundscape. The studies show that experiential complexity concerns how indoor conditions acquire significance through bodily, spatial, technological, social, and temporal relations, and through how occupants sense, interpret, anticipate, negotiate, and adjust to their surroundings over time.  This chapter answers the dissertation's research questions, brings its contributions together, and develops their implications for HBI.}

\section{research questions and dissertation contributions}
\label{sec:returning-problem}

The dissertation addressed a problem identified within HBI: human experience is central to the study of interactive buildings~\cite{alavi2019introduction,rao2023towards}, while the forms and relations encompassed by this experience require greater understanding and specification~\cite{rao2025smart} (cf.~\autoref{ch:study-one}). This section returns to the four research questions through which this problem was addressed, states the corresponding dissertation contributions, and then responds to the overarching research question.

\paragraph{\rqref{scoping}.}

\emph{What forms of human experience are addressed in research on interactive indoor environments, and how does this research seek to shape them through design and technology?}

The scoping review presented in \autoref{ch:study-one} identifies a broad
experiential agenda across research on interactive indoor environments.
Across 192 papers, it distinguishes 11 categories of human experience:
\textit{Aesthetic Experience}, \textit{Bodily Experience},
\textit{Comfort}, \textit{Democratic Experience}, \textit{Explorative
and/or Future and Speculative Experiences}, \textit{Emotional
Experience}, \textit{Experience of Relatedness}, \textit{Inclusion
Experience}, \textit{Learning and Teaching Experience}, \textit{Sense
of Orientation}, and \textit{Sense of Place-making}. This breadth is
distributed across different disciplinary perspectives, methodologies,
and design approaches.

The review further identifies 20 design mechanisms through which these
experiences are pursued. These include mechanisms such as
\textit{Spatial Design}, \textit{Community and Social Connection},
\textit{Participation and Co-design}, \textit{Somaesthetic Design},
\textit{Technological Immersion}, \textit{Personalisation}, and
\textit{Biophilic Design}. Some mechanisms support several forms of
experience, while others are more closely associated with particular
experiential aims
(cf.~\autoref{sec:ch2-findings-design-mechanisms}). The review thus
specifies both the forms of experience addressed across the field and
the different design approaches through which researchers have sought
to shape them.

Two typologies specify the technological means through which this work is undertaken. The first comprises seven intervention types: \textit{Interactive Displays}, \textit{Audio and Lighting Technologies}, \textit{Environmental Sensing}, \textit{Feedback Systems}, \textit{Augmented and Virtual Reality}, \textit{Physiological Sensing}, and \textit{Shape-Changing and Kinetic Interfaces}. The second positions technological interventions across the temporal layers of buildings by adapting \citet{brand1995buildings}'s shearing layers model and introducing \textit{Extended Reality Space} as an additional
layer. This temporal analysis shows that interactive technologies are
concentrated in the more adaptable parts of buildings, with comparatively little work engaging longer-lasting architectural layers (cf.~\autoref{sec:ch2-findings-technology}).

The first dissertation contribution is therefore a structured account of what forms of experience have been addressed, the design mechanisms through which
they have been pursued, and the technological interventions through
which these approaches have been realised. This mapping establishes the experiential landscape for the subsequent studies. It also identifies particular directions for further examination: \textit{Emotional Experience} constitutes a substantial strand of the reviewed literature, while \textit{Biophilic Design} appears as a less frequently represented design mechanism. The dissertation subsequently examines these directions in greater depth.



\paragraph{\rqref{afi}.}

\emph{What does an affective lens reveal about how occupants experience
and interpret interactive indoor environments, and how can these
insights inform HBI research and design?}


The two occupant studies presented in \autoref{ch:study-two} adopt
Affective Interaction (AfI) from HCI as the dissertation's affective
lens.\sidenote{The term \emph{Affective Interaction} was explicitly articulated in HCI by \citet{lottridge2011affective}. More recent work has used it to bring together a broader tradition of research that approaches affect as embodied, situated, and socially mediated, in contrast to approaches centred primarily on emotion recognition and classification~\cite{boehner2007emotion,ahmadpour2025affective}.} Applied to HBI, this lens brings into view how indoor conditions acquire significance through bodily sensation, activity, spatial position, social circumstances, anticipation, prior experience, and occupants' interpretations of technological behaviour.

The analysis identifies five Ways through which occupants encounter and interpret interactive indoor environments: \textit{Interpreting the Somatic Experience Over Time}, \textit{Anticipating Disruption Before It Happens}, \textit{Constructing Spatial Purpose}, \textit{Micro-Adjustments and Low-Effort Escalation}, and \textit{Making Sense of Smartness} (cf.~\autoref{sec:ch3-higher-order}). Together, the Ways show that occupants relate to indoor environments prospectively and interpretively, rather than responding only to conditions as they occur.

These findings are consolidated into six Affective Practices, each paired with an enabling Interaction Quality:
\textit{AP1---Orient and Appraise} with
\textit{IQ1---Somatic Legibility};
\textit{AP2---Select a Micro-Condition} with
\textit{IQ2---Spatial Niche Discovery};
\textit{AP3---Social Attunement and Coordination} with
\textit{IQ3---Social Cue Surfacing};
\textit{AP4---Modulate and (De)Escalate} with
\textit{IQ4---(De)Escalation-Aware Interaction};
\textit{AP5---Probe and Align with Smart Building Technology} with
\textit{IQ5---Smartness Legibility}; and
\textit{AP6---Settle and Stabilise} with
\textit{IQ6---Active Non-Intervention} (cf.~\autoref{sec:ch3-objectives}). 

The Affective Practices describe recurring ways in which occupants manage their relationship with indoor conditions, while the Interaction Qualities translate these findings into capacities that interactive environments might support.

The second dissertation contribution therefore extends Affective Interaction (AfI) into HBI as a methodological and conceptual lens for investigating situated indoor experience. It shows how affect can provide an entry point into how environmental conditions come to matter for occupants, while the Affective Practices and Interaction Qualities carry this understanding towards design. The conceptual and methodological commitments of this contribution are developed further in \autoref{sec:affect-contribution}.

\paragraph{\rqref{biophilic}.}

\emph{What possibilities does multisensory biophilic design offer for future interactive indoor environments, and how can these possibilities be articulated to support HBI design?}

The architecture-facing studies presented in \autoref{ch:study-three} and \autoref{ch:study-four} identify multisensory biophilic design as a design direction for extending interactive indoor environments. First, the expert interviews~\autoref{ch:study-three} expand the biophilic design space through eight themes: \textit{Ecological Authenticity and the Essence of Nature in Building Design}, \textit{Art as an Expression for Biophilic Experiences}, \textit{Biophilic Design for Ecological Empathy}, \textit{Multisensory and Multidimensional Biophilic Design}, \textit{Multispecies Flourishing in Buildings}, \textit{Biophilic Design for Social Justice}, \textit{The Body in Biophilic Design}, and \textit{Biophilic Design for Occupant (Dis)Comfort}. These themes expand the experiential and ecological vocabulary available for conceiving biophilic interaction in HBI.

Five design opportunities are also developed with supporting emerging technologies: \textit{Designing for Emergent Natural Processes}, \textit{Designing for Occupant Body Sensitivity}, \textit{Designing for Biophilic (Dis)Comfort}, \textit{Designing for Biophilic Equity}, and \textit{Designing for Relational Encounters with Nature}. These opportunities demonstrate how interactive indoor systems and specific technologies might support novel biophilic experiences.

The subsequent expert design sessions presented in \autoref{ch:study-four} address how such possibilities can be structured for design. Their analysis produces a three-layer Biophilic Interaction Design Framework (BIDF) comprising seven design dimensions, seven design values, and six interaction forms (cf.~\autoref{sec:ch5-overview}). The dimensions structure the range of considerations that designers and researchers can attend to when developing interactive biophilic environments. These are: \textit{Sensory Modality}, \textit{Spatial Context}, \textit{Interaction Type}, \textit{Affective Aim}, \textit{Natural Reference}, \textit{Temporal Behaviour}, and \textit{Level of Mediation}. 

The design(er) values capture the considerations that inform how these possibilities are judged and prioritised i.e. why particular configurations matter. The seven values are: \textit{Sensory Ecology}, \textit{Attentional Ease}, \textit{Proprioceptive Engagement}, \textit{Irreversibility}, \textit{Unknowability}, \textit{Non-Reciprocity}, and \textit{Collective Liability}.
 
The interaction forms describe recurring logics through which occupants might encounter and participate in biophilic indoor environments. These are: \textit{Anchoring}, \textit{Synchronising}, \textit{Dwelling}, \textit{Inferring}, \textit{Withholding}, and \textit{Accepting}. Together, the three layers make the design reasoning developed
through the expert sessions available to HBI as a structured design resource. In particular, the framework extends consideration beyond immediate responsiveness and personalisation towards interaction that may involve duration, ambiguity, limited control, environmental variation, and shared consequence.

The third dissertation contribution therefore combines an expanded
account of what multisensory biophilic interaction might encompass with
a framework for reasoning about how such interactions might be designed
at architectural scale.


\paragraph{\rqref{sound}.}

\emph{How can an AI-transformed office soundscape be developed as a form
of biophilic interaction, and what does an affective lens reveal about
how it is experienced and interpreted in shared workplace contexts?}

The sound-focused direction developed in the third expert design session
in ~\autoref{subsec:ch5-session3} considered how ongoing office sound could be
transformed into nature-like sound while preserving a perceptible
connection to the surrounding office soundscape. This direction informed
the development of the Mediated Office Soundscapes System (MOSS) \autoref{ch:study-six}. MOSS uses neural audio models
trained on nature sounds to transform incoming office audio into bird-like or water-like sound. 

This relationship between the transformed sound and ongoing office activity is conceptualised as \emph{biophilic mediation}. Biophilic interaction is therefore produced through the ongoing activity of the workplace itself. MOSS provides a working system through which this interaction design concept can be experienced and examined.

The system was examined through an online listening study and live
demonstrations with architecture and spatial-sound experts.
The listening study showed that the AI-transformed condition was
rated as more calming, less cold or clinical, less distracting, and
less out of place than conventional nature sound
(cf.~\autoref{ch:study-six}). Participants' accounts further
showed that acceptance depended on how the transformed sound related to
the office context, how strongly it occupied attention, and whether its
response to surrounding activity could be understood.

An affective lens therefore shows that the experiential value of the soundscape depends on more than the acoustic quality of the generated nature-like sound. A transformation can remain recognisably artificial while still fitting the workplace and altering its affective character. The findings identify three experiential conditions shaping biophilic mediation in shared workplaces: \emph{contextual integration}, \emph{bounded salience}, and \emph{temporal and interpretive coherence} (cf.~\autoref{subsec:moss-experiential-conditions}). Together, these conditions concern how the transformed sound relates to its context, how strongly it enters attention, and whether listeners can interpret its relationship with surrounding activity over time. The findings further identify configurability as a design consideration for responsive sound in shared environments.

The fourth dissertation contribution therefore comprises MOSS as a working instance of biophilic mediation and an account of the experiential conditions that shape how such mediation is encountered in shared workplace contexts. It shows how ongoing indoor sound can be transformed while retaining a relationship with surrounding activity and identifies \emph{contextual integration}, \emph{bounded salience}, and \emph{temporal and interpretive coherence} as conditions shaping that experience. The expert demonstrations further identify configurability as a design consideration for responsive sound in shared environments.

\paragraph{Responding to the overarching research question.}

This dissertation shows that an affective orientation enables HBI to
specify and investigate experiential complexity by examining how indoor
conditions acquire significance within situated experience. It directs
attention beyond specific environmental variables towards how
occupants sense and interpret surrounding environmental conditions, anticipate change, negotiate available possibilities, and adjust their relationship with their environment over time.

Across the dissertation, these contributions form a progression from
mapping the experiential landscape of interactive indoor environments,
to investigating how indoor conditions acquire situated significance,
developing multisensory biophilic design knowledge, and examining one
technically realised soundscape. Affective HBI connects these stages by
carrying knowledge of situated indoor experience into the development
and examination of future interactive environments.



\section{affective hbi as an orientation for situated inquiry}
\label{sec:affect-contribution}

This section develops the inquiry-oriented contribution of affective HBI. Drawing primarily on the occupant studies in~\autoref{ch:study-two}, it specifies affective HBI in terms of what it investigates, how situated experience is analysed, and how the resulting knowledge can inform design. Emotion provides an entry point into indoor experience; Affective Practices provide a unit of analysis; and situated methods retain the circumstances through which experience acquires meaning. The section then considers how these experiential relations can be carried forward as questions for future interactive environments.

Affective HBI builds on HBI's extension of interaction design to architectural scale, where buildings, spatial arrangements, environmental systems, and computational processes become part of interaction~\cite{wiberg2020interaction,schnadelbach2019adaptive}. It also draws on experience-centred HCI, which attends to experience as embodied~\cite{hook2018designing}, situated~\cite{dourish2001action}, felt, and interpretive~\cite{mccarthy2004technology,wright2008empathy}. Within HBI, research concerned with occupants' subjective and lived experience further establishes the importance of examining how people interpret and relate to technologically mediated environments~\cite{rao2023towards, margariti2023understanding}.

This position draws particularly on Affective Interaction (AfI), which treats affect as situated, embodied, socially mediated~\cite{lottridge2011affective,ahmadpour2025affective}, and constituted through interaction~\cite{boehner2005affect,boehner2007emotion}. Recent work characterises AfI as an alternative tradition within affective HCI that foregrounds the contextual and interactional constitution of affect~\cite{ahmadpour2025affective}.  Affective HBI brings this perspective into architectural-scale experience by using affect to investigate how spatial, environmental, technological, and social conditions acquire significance for occupants~\cite{paraschivoiu2022affective} (cf.~\autoref{sec:ch3-s1-findings}, \autoref{sec:ch3-s2-findings}, and \autoref{sec:ch3-disc-reflections-afi}). Emotion therefore serves as a sensitising entry point: an account of calmness, frustration, relief, or discomfort directs inquiry towards the relations and circumstances through which that experience becomes meaningful.

Within this dissertation, affective HBI is specified through three connected elements of inquiry. First, its \emph{object of inquiry} is the situated relation through which an indoor condition acquires affective significance (cf.~\autoref{sec:ch3-s1-findings} and \autoref{sec:ch3-s2-findings}). Second, its \emph{unit of analysis} is the practice through which an occupant appraises, anticipates, interprets, adjusts to, or settles within that condition (cf.~\autoref{sec:ch3-higher-order}). Treating these relations as practices keeps attention on how occupants engage with indoor conditions through sequences of interpretation and action, rather than reducing the analysis to isolated affective responses. Third, its \emph{methodological approach} retains the spatial, bodily, technological, social, and temporal circumstances through which experience is produced and interpreted (cf.~\autoref{sec:ch3-s1}, \autoref{sec:ch3-study2}, and \autoref{sec:ch3-objectives}).

These three elements are connected to design through a further step: translating identified experiential relations into questions about what future interactive environments might need to make perceptible, understandable, negotiable, adjustable, or possible. This translation retains the conditions that gave the empirical finding its significance while expressing that finding in a form that can inform subsequent design reasoning. Together, these elements specify how affective HBI produces situated knowledge about indoor experience and carries that knowledge towards the design of interactive indoor environments.


\subsection{emotion as an entry point for investigating indoor experience}
\label{subsec:emotion-entry-point}

Within the occupant studies, emotion is used as a way of investigating how environmental conditions come to matter for occupants, rather than as a state to be classified, as in approaches that computationally model affective states~\cite{picard1999affective, gao2020n}.  Reports of feeling calm, exposed, frustrated, uncertain, or relieved can direct attention to the circumstances that produced those feelings and to how occupants interpreted them~\cite{boehner2007emotion,lottridge2011affective} (cf.~\autoref{sec:ch3-s1-findings} and \autoref{sec:ch3-s2-findings}). The in-situ survey from \autoref{ch:study-two} supports this methodological role: emotion-framed accounts contained richer bodily, social, and situational information than comfort-framed accounts (cf.~\autoref{sec:ch3-s1-findings}).

The meaning of an emotion depends on the situation in which it is experienced~\cite{barrett2022context}. Calmness, for example, may relate to quietness, daylight, enclosure, freedom from interruption, or confidence that conditions will remain consistent, and these qualities matter differently when an occupant is concentrating, resting, meeting someone, or moving through the building. Likewise, quietness may support concentration in one situation while making another person's movements more conspicuous in another. Affect therefore cannot be assigned a fixed meaning to a particular environmental condition~\cite{barrett2022context,boehner2007emotion} (cf.~\autoref{sec:ch3-s1-findings} and \autoref{sec:ch3-s2-findings}).

Affective HBI uses emotion-framed inquiry to examine these relations. This complements approaches that measure~\cite{mehrabian1974approach,russell1980circumplex,bower2019impact} or model affective responses~\cite{picard1999affective,gao2020n} to environmental stimuli by asking how a condition becomes significant within situated use~\cite{boehner2005affect,boehner2007emotion,lottridge2011affective} (cf.~\autoref{sec:ch3-afi} and \autoref{sec:ch3-disc-reflections-afi}).


\subsection{affective practices as the unit of analysis}
\label{subsec:affective-practices}

Emotion provides an entry point into situated experience, while Affective Practices provide the unit of analysis. They describe the practical and interpretive work through which occupants establish, maintain, or revise their relationship with an indoor environment. This includes attending to bodily and environmental cues, anticipating change, interpreting spatial and social conditions, testing possible adjustments, and deciding when a situation is workable (cf.~\autoref{sec:ch3-s2-findings} and \autoref{sec:ch3-higher-order}).

This practice-based account shifts analytical attention from decontextualised affect ratings, such as valence--arousal ratings~\cite{russell1980circumplex,rao2023affective, kuutti2014turn}, towards how occupants act and make judgements over time. Previous encounters shape expectations, anticipated activity influences present spatial decisions, and small adjustments allow occupants to test whether a condition supports what they are trying to do. The corresponding Interaction Qualities translate these observed practices into characteristics that interactive environments might support (cf.~\autoref{sec:ch3-higher-order}).

Affective Practices also place comfort within a wider field of situated judgement. Environmental conditions~\cite{jazizadeh2014human,marques2020real,van2010effect} and reports of satisfaction~\cite{zhang2012factors,alavi2017comfort} remain important for healthy, safe, sustainable, and usable buildings, but occupants judge these conditions alongside privacy, activity, social proximity~\cite{clear2018thermokiosk,von2021towards}, familiarity, available
alternatives, and expectations of building behaviour~\cite{mitchell2020human}. A thermally imperfect location, for example, may still be preferred because it provides greater visual protection. The contribution is therefore to show how comfort forms one part of the broader relations through which occupants determine whether an environment works for them (cf.~\autoref{sec:ch2-findings-experiences}, \autoref{sec:ch2-implications-landscape}, and \autoref{sec:ch3-s2-findings}).


\subsection{situated methods for architectural-scale experience}
\label{subsec:situated-methods}
Architectural-scale experience~\cite{alavi2019introduction,becerik2022field} extends across rooms, routes, spatial arrangements~\cite{alavi2018artifacts}, technological infrastructures~\cite{schnadelbach2019adaptive}, and encounters over time~\cite{turmo2025temporal}.  Affective HBI therefore requires methods that retain the circumstances in which experience occurs~\cite{dourish2001action,mccarthy2004technology} rather than separating reported feelings from situations in which they acquire meaning~\cite{boehner2007emotion,lottridge2011affective,ahmadpour2025affective}.

The occupant studies combined emotion-framed surveys, building walks~\cite{mitchell2020no,crivellaro2015contesting}, body mapping~\cite{turmo2023towards}, situated discussion, and open-ended accounts (cf.~\autoref{sec:ch3-s1} and \autoref{sec:ch3-study2}). These methods allowed participants to locate experiences within particular parts of the building, compare spatial alternatives, demonstrate how they responded to conditions, recall previous encounters, and anticipate what might happen next. They therefore connected affective accounts to the situations in which those accounts were formed (cf.~\autoref{sec:ch3-s1-findings} and \autoref{sec:ch3-s2-findings}).

This is particularly important in interactive indoor spaces, where occupants often encounter technology indirectly~\cite{tolmie2002unremarkable,schnadelbach2019adaptive} through changes in lighting, sound, temperature, access, or other environmental conditions~\cite{saputra2023smart,daissaoui2020iot}, without necessarily knowing the sensing and control processes behind them~\cite{bellotti2002making,lim2009and} (cf.~\autoref{sec:ch3-s2-findings} and \autoref{sec:ch3-higher-order}). Situated methods make it possible to examine how occupants interpret such behaviour~\cite{bellotti2002making,lim2009and} and how those
interpretations influence subsequent action~\cite{dourish2001action,hu2023scoping}. Their methodological contribution to HBI is therefore to provide explanatory depth while retaining the architectural circumstances through which experience acquires meaning (cf.~\autoref{sec:ch3-objectives}, \autoref{sec:ch3-discussion}, and \autoref{sec:ch3-disc-reflections-afi}).


\subsection{from experiential findings to design questions}
\label{subsec:experiential-design-knowledge}

The preceding elements specify how affective HBI investigates situated
indoor experience. Their relevance to design lies in carrying the
relations identified through that inquiry forward as design questions.
Affective inquiry can support multiple design responses to the same experiential relation.
An account of feeling exposed, for example, may reveal a relation among
visibility, activity, social presence, and available spatial
alternatives. The resulting design question concerns how those
relations might be supported, altered, made legible, or deliberately
left open, rather than how the reported emotion itself should be
removed~\cite{rao2025lessons}.

The Affective Practices and Interaction Qualities provide one
illustration of this translation
(cf.~\autoref{sec:ch3-higher-order}). Practices describe how occupants
orient, interpret, negotiate, adjust, and settle within environmental
conditions; the corresponding Interaction Qualities identify
capacities that an interactive environment might support. They
therefore preserve the relation identified through inquiry while
shifting the question from what an occupant experienced to what a
future environment might enable.


\section{architecture as design knowledge for affective hbi}
\label{sec:architecture-design-knowledge}

Architectural knowledge provides the next step in developing possible
responses to the design questions produced through affective inquiry. HBI has called for interaction design to engage buildings as persistent, shared, and spatially distributed settings of computation~\cite{alavi2018artifacts,alavi2019introduction,wiberg2020interaction} and to work with the practitioners and domain experts who shape them~\cite{alavi2019introduction,kirsh2019architects,rao2025smart}.The dissertation takes up this direction through sustained engagement with architectural and spatial expertise across expert interviews~\cite{rao2025designing} (cf.~\autoref{sec:ch4-methodology}), expert design sessions~\cite{rao2026biophilic} (cf.~\autoref{subsec:ch5-study-design}), and live prototype demonstrations~\cite{rao2027biophilic} (cf.~\autoref{ch:study-six}).

Architecture contributes here as both the setting of interaction and a source of design knowledge. Architectural expertise brings ways of reasoning about how spatial, sensory, material, environmental, and temporal conditions are composed and inhabited~\cite{pallasmaa2024eyes,griffero2016atmospheres,zumthor2006atmospheres,verma2022appraisal}. 

The architecture-facing studies bring this knowledge into conversation with embodied, experience-centred~\cite{hook2018designing}, soma-based~\cite{tschaepe2021somaesthetics}, and material~\cite{van2018skin} and spatial approaches in HCI~\cite{nabil2017interioractive}. This exchange broadens the forms of knowledge through which HBI can conceive and judge future interactive environments~\cite{rao2025designing,rao2026biophilic}.

The section develops this contribution in two steps. It first examines how architectural expertise can function as a design resource within HBI, and then considers how this exchange changes the object of interaction design at architectural scale.


\subsection{architectural expertise as a design resource}
\label{subsec:architectural-expertise}

Engaging architectural expertise provides access to forms of domain knowledge developed through architectural practice and creates ways of bringing that knowledge into interaction-design inquiry~\cite{kirsh2019architects}. In the expert interviews~\cite{rao2025designing} (cf.~\autoref{sec:ch4-methodology}), architectural and spatial practitioners articulated knowledge, experience, and future possibilities that expanded the design questions available to HBI. In the subsequent design sessions~\cite{rao2026biophilic} (cf.~\autoref{subsec:ch5-study-design}), HCI design methods provided a means of externalising and examining this knowledge through concrete design proposals.

The Biophilic Interaction Design Framework~\cite{rao2026biophilic} (cf.~\autoref{sec:ch5-overview}) is one outcome of this exchange. Its dissertation-level significance lies in demonstrating that forms of architectural practitioner knowledge can be elicited and structured as a design resource for HCI. This process makes architectural design judgements available for analysis, comparison, and further development within interaction design.

The subsequent live demonstrations of MOSS~\cite{rao2027biophilic} (cf.~\autoref{ch:study-six}) extended this exchange to an emerging form of AI-based interaction. They asked architectural and spatial-sound experts to consider
computational behaviour itself as a matter of spatial design within
the architectural setting. \sidenote{This boundary is already beginning to shift in architectural practice. In discussions surrounding the studies, some practitioners described recent projects for AI technology-based companies in which interactive lighting and other responsive systems formed part of workplace design. Such examples suggest increasing architectural engagement with computationally mediated environments, although these forms of interaction are not yet routine across practice.}

Across these engagements, knowledge moved in both directions~\cite{rao2025designing,rao2026biophilic,rao2027biophilic}. Architectural practitioners brought domain knowledge from architectural practice into HCI research, while HCI methods provided ways of eliciting, structuring, technically developing, and critically examining that knowledge.  The collaboration provides a means through which architectural and interaction-design knowledge can shape one another in the development of interactive indoor environments. This offers one response to HBI's call for closer engagement with the disciplines responsible for the design of built environments~\cite{alavi2018artifacts,alavi2019introduction,kirsh2019architects}.


\subsection{architectural knowledge and the scope of interaction design}
\label{subsec:architectural-inhabitable-relation}

Architectural expertise broadens what interaction design must consider at building scale~\cite{kirsh2019architects}. It directs attention from computational behaviour itself towards the spatial and environmental conditions affected by its operation~\cite{nguyen2022responsive}. A computational change may alter conditions beyond the person who initiated it~\cite{niemantsverdriet2018share}, remain perceptible after the initiating action has passed, or acquire different significance as occupants move, activities change, and other people enter the setting~\cite{alavi2018artifacts,nguyen2022responsive,nguyen2024adaptive}. Studies of shared environmental systems similarly show that interaction with environmental controls can involve coordination among multiple occupants~\cite{niemantsverdriet2018share,vanDeWerff2019lighting}.

Timing, spatial distribution~\cite{nguyen2022responsive,nguyen2024adaptive}, and consequences for other occupants~\cite{niemantsverdriet2018share,vanDeWerff2019lighting} therefore become matters of interaction-design judgement alongside system behaviour. Architectural knowledge contributes ways of considering how computational behaviour becomes part of conditions that are spatially distributed, persistent, and shared~\cite{alavi2018artifacts,alavi2019introduction,wiberg2020interaction}.

These concerns are particularly relevant to responsive environments, where occupants may need to interpret what a system is doing~\cite{bellotti2002making,lim2009and} and where automated behaviour coexists with the activities and preferences of multiple occupants~\cite{vanDeWerff2019lighting,niemantsverdriet2018share}. The following section examines these relations through the technically realised case of MOSS.



\section{moss as a design instance}
\label{sec:sound-instance}

MOSS provides a concrete case for examining how computational
behaviour becomes part of the environmental conditions occupants encounter. It technically develops one sound-focused biophilic direction from the expert design sessions~\cite{rao2026biophilic} (cf.~\autoref{subsec:ch5-session3}) into an AI-transformed office soundscape~\cite{rao2027biophilic}.
The affective orientation developed in \autoref{sec:affect-contribution}, drawing on Affective Interaction~\cite{boehner2005affect,boehner2007emotion,lottridge2011affective}, shapes this examination. Attention is directed towards how listeners relate transformed sound to surrounding activity, how system behaviour enters attention, and whether the resulting condition fits their activity and workplace situation~\cite{deacon2024working} (cf.~\autoref{ch:study-six}). MOSS thereby brings the dissertation's biophilic design work and affective inquiry together in a technically realised system.

\subsection{biophilic mediation as an interaction design concept}
\label{subsec:biophilic-mediation}

MOSS develops \emph{biophilic mediation} as an interaction design concept~\cite{rao2027biophilic}. I use \emph{mediation} to describe an altered relationship between a listener and an existing environmental condition. In MOSS, sounds occurring in the indoor environment are transformed into nature-like sound while characteristics of the incoming audio continue to shape what the listener hears (cf.~\autoref{ch:study-six}). The system therefore changes how ongoing acoustic activity is encountered while retaining a perceptible, variable, and partly ambiguous connection to that activity~\cite{gaver2003ambiguity,bustoni2024exploring}.

This distinguishes biophilic mediation from conventional nature playback~\cite{newbold2017nature,yu2016auditory}, where an independent nature recording is introduced into the soundscape, and from acoustic masking, which commonly seeks to reduce the prominence of unwanted sound~\cite{haas2018soundscape,jacobsen2024personal, bustoni2024exploring}.

Biophilic mediation was developed around three \emph{design aims}: \emph{enriching quietness} through low-level nature-like variation, \emph{reframing activity} by changing the auditory character of office events, and \emph{preserving environmental connection} by allowing changes in surrounding activity to remain consequential to the soundscape~\cite{rao2027biophilic} (cf.~\autoref{ch:study-six}). These aims define the intended form of interaction. The following subsection identifies the experiential conditions that shaped how listeners encountered it.


\subsection{experiential conditions of biophilic mediation}
\label{subsec:moss-experiential-conditions}

The studies identify three connected experiential conditions that shape how biophilic mediation is encountered in an office context: \emph{contextual integration}, \emph{bounded salience}, and \emph{temporal and interpretive coherence}~\cite{rao2027biophilic} (cf.~\autoref{ch:study-six}). These conditions connect the preceding design aims to how listeners experienced and interpreted the realised soundscape.

The first condition is \emph{contextual integration}. In the listening study, the AI-transformed sound was rated as less nature-like than conventional nature sound while receiving comparable pleasantness ratings, higher calmness ratings, and lower ratings for cold or clinical character, distraction, and being out of place~\cite{rao2027biophilic} (cf.~\autoref{ch:study-six}). Participants also recognised the intended natural referents while reporting artificiality, artefacts, restricted variation, and discontinuity (cf.~\autoref{ch:study-six}).

These findings distinguish natural realism from contextual integration as two qualities of the listening experience. A recognisably artificial transformation can establish a coherent relationship with office activity, while an independent nature recording can remain experientially detached from that activity. \emph{Contextual integration} concerns how the transformed sound relates to the spatial, temporal, social, and activity conditions in which it is encountered~\cite{thi2026envisioning}. For biophilic mediation, this relationship with the surrounding environment constitutes an experiential quality alongside the recognisability and realism of the natural reference~\cite{newbold2017nature,yu2016auditory,riaz2025bird}.

The second condition is \emph{bounded salience}. The transformed sound remains perceptible while occupying a limited position in focal attention. This concern relates to work on calm and unremarkable interaction~\cite{case2015calm,tolmie2002unremarkable,harman2016notifall}. In the listening study, soft, continuous, and proportionate output could alter the felt character of the office while remaining peripheral, whereas sharp events, abrupt transitions, repetition, and excessive sensitivity increased attentional demand~\cite{rao2027biophilic} (cf.~\autoref{ch:study-six}).


Within the dissertation, bounded salience provides a more specific account of the \emph{Attentional Ease} developed in the Biophilic Interaction Design Framework~\cite{rao2026biophilic} (cf.~\autoref{sec:ch5-layer2}). It concerns how strongly and how frequently computational behaviour becomes perceptually prominent over time. The soundscape can contribute atmospheric variation and retain a perceptible relationship with activity while allowing occupants' attention to remain directed towards the activities for which the setting is being used.

The third condition is \emph{temporal and interpretive coherence}. Participants differed in whether they perceived the system as responsive and in the relationships they inferred between office activity and transformed output~\cite{rao2027biophilic} (cf.~\autoref{ch:study-six}). The experience of responsiveness depends on whether listeners can relate changes in the soundscape to preceding events and form a workable account of the system's broad behaviour \cite{bellotti2002making,lim2009and}. In MOSS, coherent relationships between activity and output could support liveliness and environmental connection, whereas highly salient, delayed, inconsistent, or difficult-to-anticipate changes could amplify activity, expose the mechanism, or increase self-consciousness (cf.~\autoref{ch:study-six}). \emph{Temporal and interpretive coherence} therefore concerns whether computational change remains sufficiently connected to ongoing activity for listeners to interpret its behaviour over time.

The expert demonstrations place these experiential conditions in relation to spatial implementation~\cite{rao2027biophilic} (cf.~\autoref{ch:study-six}). Experts considered personal, localised, room-level, and building-integrated forms of responsive sound and related their suitability to room purpose, acoustic character, activity, duration, and delivery format (cf.~\autoref{ch:study-six}).  Their requests for influence over sensitivity, timing, auditory character, intensity, and system operation indicate how these conditions may need to be considered when configuring responsive sound in shared environments (cf.~\autoref{ch:study-six}). At architectural scale, their realisation would depend on where the intervention is encountered, how its behaviour is distributed through space and time, and what opportunities occupants are given to influence that behaviour.

Together, the three conditions locate the experiential contribution of biophilic mediation in the relationships among computational behaviour, ongoing environmental activity, attention, interpretation, and use. They specify how an AI-transformed soundscape becomes an experiential condition of the workplace and provide the basis for examining how the technical properties of MOSS shape that condition.


\subsection{technical properties as environmental qualities}
\label{subsec:technical-experiential}

MOSS shows how technical properties and implementation decisions become environmental qualities when computational behaviour shapes a condition that occupants encounter~\cite{alavi2018artifacts,alavi2019introduction,wiberg2020interaction}. Latency, smoothing, model behaviour, sensing configuration, mixing, and playback configuration shape perceptible qualities of the soundscape, including its continuity, prominence, timing, and the extent to which ongoing office activity remains audible (cf.~\autoref{ch:study-six}).


In MOSS, latency affects whether an auditory change appears related to an office event; smoothing shapes whether the output is experienced as a continuous texture or as a sequence of discrete reactions; and model artefacts can increase attentional demand and alter the perceived naturalness of the sound (cf.~\autoref{ch:study-six}). Microphone placement determines which activities become consequential to the transformation, while playback level and mixing shape the prominence of the transformed sound and the degree to which ongoing office activity remains perceptible (cf.~\autoref{ch:study-six}). These technical decisions therefore participate directly in producing the experiential conditions described in the preceding subsection.

For MOSS, technical evaluation therefore includes how implementation choices alter the auditory conditions through which the system is experienced. The microphone input, preprocessing, neural model, temporal conditioning, mixing, playback, room acoustics, surrounding activity, and listening situation together shape the resulting soundscape (cf.~\autoref{ch:study-six}). Evaluating the system consequently requires relating its computational behaviour to the sensory and situational conditions produced through that behaviour.

The contribution here is specific to the sound-based case examined through MOSS. OOther environmental interventions, including lighting, airflow, temperature, projection, and kinetic systems~\cite{rao2025smart} (cf.~\autoref{sec:ch2-findings-technology}), involve different perceptual, safety, spatial, and infrastructural considerations. MOSS nonetheless demonstrates how implementation decisions can become part of the environmental conditions through which a computational system is experienced.



\section{interaction at building scale}
\label{sec:interaction-building-scale}

The preceding discussion provides a basis for considering what an affective orientation implies for interaction at the scale of indoor environments. This section examines what constitutes interaction at building scale, the role occupants take within it, and how responsiveness can be understood within this account.


\subsection{what constitutes interaction at building scale}
\label{subsec:interaction-environmental-relation}

The scoping review distinguishes the technological forms introduced into buildings from the design mechanisms through which particular experiences are pursued~\cite{rao2025smart} (cf.~\autoref{sec:ch2-findings-design-mechanisms} and \autoref{sec:ch2-findings-technology}). The occupant studies show why this distinction matters in situated use. Environmental and technological conditions acquire significance through their relation to bodily state, activity, spatial position, social circumstances, prior encounters, expectations, and interpretations of building behaviour~\cite{rao2026emotion} (cf.~\autoref{sec:ch3-higher-order}). These findings show that technological conditions acquire significance through the situations in which their effects are encountered and interpreted.

The architecture-facing studies extend this relation into design. The expert interviews broaden the sensory, bodily, spatial, temporal, social, and ecological possibilities considered for interaction~\cite{rao2025designing} (cf.~\autoref{sec:ch4-design-opportunities}), while the Biophilic Interaction Design Framework structures how such possibilities can be configured through the environment~\cite{rao2026biophilic} (cf.~\autoref{sec:ch5-overview} and \autoref{sec:ch5-layer3}). Together, these studies place interaction within the spatial and environmental conditions through which occupants encounter technological behaviour~\cite{alavi2018artifacts,alavi2019introduction,wiberg2020interaction}.

MOSS makes this relationship technically concrete. Its implementation choices affect the auditory conditions produced by the system and how listeners relate those conditions to surrounding activity~\cite{rao2027biophilic} (cf.~\autoref{ch:study-six}). At building scale, technological behaviour therefore becomes part of interaction through the spatial and environmental conditions in which occupants encounter its consequences~\cite{alavi2018artifacts,nguyen2022responsive}.



\subsection{the occupant's role in interaction}
\label{subsec:occupant-participation-interaction}
The occupant studies provide a situated account of the occupant's role within these relations. People appraised conditions, anticipated change, selected spatial alternatives, coordinated with others, tested interpretations of building behaviour, made small adjustments, and sometimes retained an existing condition once a workable relation had been established~\cite{rao2026emotion} (cf.~\autoref{sec:ch3-higher-order}). Occupant involvement therefore included interpretation, spatial and behavioural adjustment, and direct changes to environmental conditions.

The biophilic studies extend this account towards interactions in which occupant involvement does not depend on continuous control. The expert studies consider bodily and social involvement alongside environmental processes that may remain partly outside individual control~\cite{rao2025designing} (cf.~\autoref{sec:ch4-design-opportunities}). The subsequent framework develops forms of interaction involving sustained engagement, interpretation, non-response, and accommodation of conditions that cannot always be individually adjusted~\cite{rao2026biophilic} (cf.~\autoref{sec:ch5-layer3}). These studies therefore broaden the occupant's role beyond operating a system towards establishing and maintaining a relationship with changing conditions.

MOSS adds direct influence over system behaviour as a design consideration.In the live demonstrations, experts requested ways to influence sensitivity, timing, sound character, intensity, and whether the system operated~\cite{rao2027biophilic} (cf.~\autoref{ch:study-six}). These requests identify possible forms of configuration rather than established preferences for a deployed shared system.

Across the dissertation, the occupant's role therefore ranges from interpreting and anticipating environmental behaviour, through spatial and behavioural adjustment and accommodation, to direct influence and coordination with others. Occupants can alter the environment, alter their own relationship with it, or negotiate conditions that are also experienced by other people.


\subsection{responsiveness within interaction}
\label{subsec:responsiveness-form-interaction}

Within this account, responsiveness is one form of interaction in which
environmental or technological behaviour changes in relation to
occupant activity, sensed conditions, explicit input, or other system
states~\cite{vallgaarda2014dress,nguyen2022responsive}. The scoping
review identifies technical interventions through which such change can
be produced, including environmental sensing, feedback systems, and
shape-changing or kinetic interfaces~\cite{rao2025smart}
(cf.~\autoref{sec:ch2-findings-technology}).

The dissertation shows that the experience of responsiveness also
depends on how environmental change is interpreted across time.
Occupants in the situated studies anticipated changes in environmental
and automated conditions, tested interpretations of building behaviour,
and revised their expectations through subsequent encounters
~\cite{rao2026emotion}
(cf.~\autoref{sec:ch3-higher-order}). The Biophilic Interaction Design
Framework further considers temporal patterns, unpredictability,
irreversibility, non-reciprocity, and forms of interaction in which an
environment does not immediately answer an occupant's action
~\cite{rao2026biophilic}
(cf.~\autoref{sec:ch5-layer1},
\autoref{sec:ch5-layer2}, and
\autoref{sec:ch5-layer3}). Responsiveness can therefore be experienced
across successive encounters rather than through a discrete
action--response exchange.

MOSS makes this distinction concrete. Incoming office activity shapes
the transformed sound~\cite{rao2027biophilic}
(cf.~\autoref{ch:study-six}), yet listeners varied in whether they
perceived this relationship and how they interpreted the resulting
changes. The experiential conditions identified through MOSS
(cf.~\autoref{subsec:moss-experiential-conditions}) show that a
system's capacity to respond does not by itself determine how that
response is experienced.

Responsiveness is therefore one way technological behaviour can become
part of building-scale interaction. Its experiential quality depends
on how environmental change relates to ongoing activity, how occupants
interpret that relationship across time, and what possibilities they
have for adjusting to or influencing its consequences.


\section{reflexivity}
\label{sec:reflexivity}

The researcher occupied different roles across the programme,
including reviewer, field researcher, interviewer, workshop
facilitator, analyst, and prototype designer. These roles shaped both
the questions asked and the material produced. Emotion-framed prompts
directed attention towards affective experience; biophilic framing
oriented expert discussion towards sensory and ecological
possibilities; and describing MOSS as AI-transformed and nature-like
provided interpretive context for participants.

The dissertation therefore maintains the provenance of its different
claims. The occupant constructs derive from situated accounts; the
biophilic themes and framework derive from expert interviews and design
sessions; and the MOSS findings distinguish comparative listening
evidence from expert critique and projected applications. The
cross-study relationships developed in this chapter are
dissertation-level interpretations across studies undertaken for
different purposes rather than findings from a single comparative
study.

\section{limitations and future work}
\label{sec:limitations-future}

This section considers the principal limitations of the dissertation and the research needed to examine the transferability and further development of its contributions.


\subsection{affective hbi across settings}
\label{subsec:limitations-affective}

The principal occupant studies centred on LAB42%
\sidenote{LAB42 is a recently built university building with flexible
shared spaces, integrated environmental sensing and automation, and
limited direct occupant control, making it a useful setting for studying
everyday experience in interactive indoor environments.}
The Affective Practices and Interaction Qualities are therefore most
directly grounded in comparable institutional and workplace settings.
Other building types, populations, cultures, and forms of building
technology may involve different forms of appraisal, authority,
adjustment, social coordination, and environmental control. The studies
also provide limited evidence about seasonal change, long-term
adaptation, and sustained relationships with newly introduced systems.

Further comparative and longitudinal research should examine which
Affective Practices and Interaction Qualities recur across settings and
how they vary with sensory needs, available alternatives, environmental
control, organisational conditions, and duration of occupation. This
work can combine situated qualitative methods with environmental,
behavioural, and physiological data while retaining the contextual and
interpretive account of experience developed in this dissertation.


\subsection{multisensory biophilic design in occupied environments}
\label{subsec:limitations-biophilic}

The expert interviews and design sessions provide professional accounts,
future-oriented possibilities, and design reasoning rather than evidence
of occupant needs or universal design principles. Most participating
experts worked within European professional contexts, and the resulting
knowledge may therefore reflect particular assumptions about climate,
sensory preference, seasonality, professional practice, and relations
with nature.

The Biophilic Interaction Design Framework should therefore be examined
through the development and use of concrete interventions in occupied
environments. Future work should investigate how its dimensions, design
values, and interaction forms remain useful when considered alongside
accessibility, engineering, regulation, maintenance, climate,
organisational practice, and everyday use. It should also involve
occupants and other affected stakeholders directly and examine the
framework across different sensory modalities and multisensory
combinations.


\subsection{moss in sustained workplace use}
\label{subsec:limitations-sound}

MOSS is an early-stage, headphone-based research prototype. The online
study examines immediate responses to controlled sound conditions, while
the live demonstrations provide professional critique of a working
relationship between office activity and transformed sound. Study labels
and the nature-oriented framing may also have influenced how participants
interpreted the sound. The studies therefore do not establish
habituation, long-term workplace acceptance, room-level collective
experience, or continuous technical operation.

The next step is deployment in shared workplaces over longer periods.
Such studies should examine how \emph{contextual integration},
\emph{bounded salience}, and \emph{temporal and interpretive coherence}
are affected by repeated exposure, room acoustics, changing activity,
habituation, technical artefacts, and different forms of configuration.
Deployment would also allow investigation of how several occupants
interpret and influence the same soundscape, how preferences are
negotiated, and how privacy, responsibility, and control are handled in
everyday workplace use.

\section{conclusion}
\label{sec:final-conclusion}

Human experience in interactive indoor environments is shaped through
relations among environmental conditions, technological behaviour,
activities, social situations, and occupants' interpretations and
actions. This dissertation has shown how an affective orientation can
investigate these relations and carry the resulting knowledge into
the design of future interactive indoor environments.

Across the four contributions, the dissertation moves from specifying
the experiential landscape of interactive indoor environments, to
investigating situated experience through affect, developing
multisensory biophilic design knowledge, and examining one
AI-transformed soundscape as a concrete design instance. Together, these
studies show how experiential knowledge can be carried from inquiry into
the development and examination of future interactive environments.

For HBI, interaction at building scale therefore concerns how
technological behaviour becomes part of the environmental conditions
occupants encounter, how occupants interpret and adjust to those
conditions, and how they may influence their consequences. Together,
these studies position affective HBI as a way to understand, imagine,
and design indoor environments from the complexity of human experience.






